% !LW recipe = latexmk (xelatex)
% 双栏终局笔记。LaTeX Workshop 使用上面的 XeLaTeX 配方。
\documentclass[UTF8,a4paper,10pt]{ctexart}
\usepackage[margin=14mm,top=13mm,bottom=15mm]{geometry}
\setlength{\columnsep}{5mm}
\setlength{\parindent}{1em}
\setlength{\parskip}{0pt}
\linespread{0.96}
\ctexset{
  section={format=\raggedright\bfseries\normalsize,beforeskip=7pt,afterskip=3pt},
  subsection={format=\raggedright\bfseries\normalsize,beforeskip=5pt,afterskip=2pt}
}
\usepackage{microtype}
\usepackage{xcolor}
\definecolor{attackRed}{HTML}{C03736}
\definecolor{threatYellow}{HTML}{E5B536}
\definecolor{survivalGreen}{HTML}{54A66A}
\definecolor{moveBlue}{HTML}{245E91}
\usepackage{chessboard}
\setchessboard{boardfontsize=17pt,showmover=false}

% chessboard 的标记继承棋子字体；箭头数字必须显式切换到正文字体。
\newfontfamily\arrowlabelfont{Latin Modern Sans}
% 数字大小改 \fontsize 的第一个值；白底/蓝底大小改 \fboxsep。
\newcommand{\ArrowLabel}[3]{%
  \begingroup\setlength{\fboxsep}{0.5pt}%
  \colorbox{#1}{\textcolor{#2}{{\arrowlabelfont\fontsize{6}{6}\selectfont #3}}}%
  \endgroup}
\cbDefinePgfMoveStyle{stepone}{%
  \pgfsetstrokecolor{moveBlue}%
  \pgfsetlinewidth{1.1pt}%
  \pgfsetarrowsend{stealth}%
  \pgfpathmoveto{\pgfpointorigin}%
  \pgfpathlineto{\pgfpointanchor{#1}{center}}%
  \pgfusepath{stroke}%
  \pgftext[at=\pgfpointscale{0.5}{\pgfpointanchor{#1}{center}}]{%
    \ArrowLabel{moveBlue}{white}{1}}}
\cbDefinePgfMoveStyle{steptwo}{%
  \pgfsetstrokecolor{moveBlue}%
  \pgfsetlinewidth{1.1pt}%
  \pgfsetarrowsend{stealth}%
  \pgfpathmoveto{\pgfpointorigin}%
  \pgfpathlineto{\pgfpointanchor{#1}{center}}%
  \pgfusepath{stroke}%
  \pgftext[at=\pgfpointscale{0.5}{\pgfpointanchor{#1}{center}}]{%
    \ArrowLabel{moveBlue}{white}{2}}}
% #2 is the FEN; #3 is the move sequence printed under the board.
\newcommand{\Position}[3][]{%
  \par\noindent\makebox[\columnwidth][c]{\chessboard[setfen={#2},#1]}\par\nobreak
  \noindent\makebox[\columnwidth][c]{{\footnotesize\bfseries #3}}\par%
}

\begin{document}
\twocolumn[{%
  \begin{center}
    % 少子杀王定式
    {\Large\bfseries Minor-Piece Checkmate Patterns\par}
    \vspace{2pt}
    % 双象与马象
    {\small Two Bishops and Bishop and Knight \quad\textbar\quad
      \the\year/\the\month/\the\day\par}
  \end{center}
  \vspace{-7pt}
}]

% 双象杀王
\section{BB: Two Bishops}

% 逼近：王跟进，双象封线
\subsection{Approach: King Opposition \& Bishop Block}

% 对王时，以象插入双王之间，迫使黑王退向底线；然后继续对王。
When the kings are in opposition, place a bishop between them to force the black king toward the back rank; then regain the opposition.

\Position[
  backstyle=color,color=attackRed!42,backfields={d6,e6,f6},
  color=threatYellow!55,backfields={c8,d7,f7,g8},
  color=survivalGreen!55,backfields={d8,e8,f8},
  markstyle=stepone,markmove=d5-e6
]{8/4k3/8/3BK3/5B2/8/8/8 w - - 0 1}{1.Be6}

% 常见情况：黑王g8
\subsection{Common Case: Black King on g8}

% 黑格象可在 a3--f8--h6 斜线上的任一安全位置，双象先后踩住g8和h8。图中 \textbf{1.Be6+ Kh8 2.Bf6\#}。
The dark-squared bishop may occupy any safe square along the a3--f8--h6 diagonal. The two bishops successively cover g8 and h8.

\Position[
  backstyle=color,color=attackRed!42,backfields={f7,g7,h7},
  color=attackRed!55,backfields={f8},
  color=survivalGreen!55,backfields={h8},
  markstyle=border,color=threatYellow,markfields={f7,g7,g8,h8},
  markstyle=stepone,markmove=d7-e6,
  markstyle=steptwo,markmove=e7-f6
]{6k1/3BB3/6K1/8/8/8/8/8 w - - 0 1}{1.Be6+ Kh8 2.Bf6\#}

% 罕见情况：黑王h8
\subsection{Rare Case: Black King on h8}
% 黑格象只能在 h6（理论上可以在f8，但实战应该走不出来），双象先后踩住h8和g8。图中\textbf{1.Bg7+ Kg8 2.Be6\#}。
\textcolor{attackRed}{The dark-squared bishop can only be on h6} (theoretically also f8, though that position should not arise in practice). The two bishops successively cover h8 and g8.

\Position[
  backstyle=color,color=attackRed!42,backfields={f7,g7,h7},
  color=attackRed!55,backfields={f8},
  color=survivalGreen!55,backfields={g8},
  markstyle=border,color=threatYellow,markfields={f7,g8,h8},
  markstyle=stepone,markmove=h6-g7,
  markstyle=steptwo,markmove=d7-e6
]{7k/3B4/6KB/8/8/8/8/8 w - - 0 1}{1.Bg7+ Kg8 2.Be6\#}

% 马象杀王
\section{Checkmating with Bishop and Knight}

\end{document}
